\chapter{Worker-Loss Instrumentation Baked Into the Metrics (Phase 1.4)}
\label{chap:p1-4}

\section{Goal}

Three of the six pilot metrics --- descriptive accuracy, robustness, and bounded completeness --- work
by re-running the safety proxy on a modified image (a mask or a perturbation) and checking whether the
answer flips. The pilot discovered, and documented across Chapters~5, 9, and 10, that a specific
artifact contaminates all three: when the modification makes the \emph{worker} undetectable, the proxy
falls out of its per-worker check (\texttt{if not worker\_boxes}) into a cruder scene-level fallback,
and the answer can flip for a reason that has nothing to do with the safety concept being tested. A
flip caused this way is a false signal of ``the region/answer was sensitive to the safety cue.''

The pilot handled this reactively and expensively. Chapter~10, in particular, had to run 159 fresh
top-1-mask GPU reruns \emph{solely} to recover the post-mask worker boxes that Chapter~5's descriptive
CSV never logged, just to estimate how many ``supported'' completeness verdicts were really
worker-loss artifacts. The scale-up plan (item 3.4) prescribes making this a standard, built-in part
of every affected metric's output: record post-operation worker boxes, emit a
\texttt{flip\_due\_to\_worker\_loss} flag, and report each metric twice --- \textbf{raw} (all flips) and
\textbf{genuine} (worker-loss flips excluded).

\section{What Was Done}

\subsection{A worker-loss flag in every affected metric}

Each metric's pure comparison core now derives the flag from information it already had access to but
was discarding:

\begin{itemize}
  \item \textbf{Descriptive accuracy} gained a pure core, \texttt{\_descriptive\_from\_results}, that
  computes \texttt{worker\_lost\_top1/top2} (baseline had a worker, masked rerun has none) and
  \texttt{flip\_due\_to\_worker\_loss\_top1/top2} (an answer flip that co-occurs with worker loss).
  Extracting the pure core also made the logic unit-testable without a GPU, which it previously was
  not.
  \item \textbf{Robustness} already exposed \texttt{worker\_lost}; it gained the derived
  \texttt{flip\_due\_to\_worker\_loss = answer\_changed AND worker\_lost}, so the correction is a first-
  class metric field rather than something each analysis must recompute.
  \item \textbf{Completeness} gained \texttt{classify\_sample\_worker\_loss\_corrected}, which counts a
  sample ``supported'' only if a \emph{genuine} flip (not attributable to worker loss) remains at
  top-1 or top-2.
\end{itemize}

Twelve unit tests were written test-first across the three metrics, covering the flip/no-flip $\times$
worker-lost/worker-present matrix and the corrected-verdict downgrade logic.

\subsection{Reporting, gated and record-safe}

A \texttt{config.REPORT\_WORKER\_LOSS\_CORRECTED} flag (default False) threads through scripts 05, 09,
and 10 as \texttt{-{}-report-worker-loss-corrected}. Default off reproduces the frozen metric CSVs
byte-for-byte; on, each script logs the worker-loss columns, reports raw and genuine rates side by
side, and writes to a \texttt{\_wlc}-suffixed file so no frozen artifact is overwritten.

The most consequential structural change is in completeness. Because descriptive accuracy now logs the
post-mask worker boxes itself, the bounded-completeness verdict becomes a \emph{pure rollup} of the
\texttt{\_wlc} descriptive CSV --- the 159 fresh GPU reruns Chapter~10 needed are simply gone,
replaced by reading two columns that already exist.

\section{Results}

\subsection{Descriptive accuracy: raw vs.\ genuine}

\begin{center}
\begin{tabular}{p{4.6cm}p{2.6cm}p{2.6cm}p{3.0cm}}
\toprule
\textbf{ranking} & \textbf{top-1 raw} & \textbf{top-1 genuine} & \textbf{worker-loss flips (top-1)} \\
\midrule
area (frozen)  & 36.2\% & 34.4\% & 3 of 59 \\
rule\_aware    & 38.7\% & 37.4\% & 2 of 63 \\
\bottomrule
\end{tabular}
\captionof{table}{Descriptive accuracy, raw vs.\ worker-loss-corrected. The area raw number reproduces
the frozen pilot's 36.2\% exactly. Under rule\_aware ranking the top-1 worker-loss count falls (3
$\rightarrow$ 2) because the worker is no longer the top-1 masked region --- confirming, from the
opposite direction, the Phase~1.1 observation that the artifact relocates to the top-2 stage under the
fixed ranking (top-2 genuine is 31.3\% under both rankings).}
\label{tab:p14-da}
\end{center}

\subsection{Robustness: occlusion is where worker loss bites}

Computed directly from the frozen \texttt{robustness.csv} (which already carried
\texttt{answer\_changed} and \texttt{worker\_lost}), the correction is negligible for most
perturbations and material for exactly the one that physically covers the worker:

\begin{center}
\begin{tabular}{p{4.0cm}p{2.6cm}p{2.6cm}p{3.4cm}}
\toprule
\textbf{perturbation} & \textbf{raw} & \textbf{genuine} & \textbf{worker-loss flips} \\
\midrule
occlude          & 28.2\% & \textbf{22.7\%} & 9 of 163 \\
low\_light       & 11.7\% & 11.0\% & 1 of 163 \\
blur             & 27.0\% & 27.0\% & 0 \\
contrast\_shift  & 10.4\% & 10.4\% & 0 \\
reworded\_prompt & 27.5\% & 27.5\% & 0 \\
\midrule
overall          & 20.8\% & 19.5\% & 10 of 790 \\
\bottomrule
\end{tabular}
\captionof{table}{Robustness answer-change rate, raw vs.\ genuine. Occlusion's 5.5-point inflation is
concentrated worker loss --- the centered occlusion square covers the worker and reroutes the proxy ---
directly supporting the plan's separate call (Phase~2.4) to split targeted from random occlusion.}
\label{tab:p14-rob}
\end{center}

\subsection{Completeness: the 159 reruns, replaced by a rollup}

\begin{center}
\begin{tabular}{p{4.4cm}p{3.2cm}p{3.2cm}p{2.4cm}}
\toprule
\textbf{ranking} & \textbf{supported (raw)} & \textbf{supported (corrected)} & \textbf{downgraded} \\
\midrule
area (frozen)  & 71/163 (43.6\%) & 67/163 (41.1\%) & 4 \\
rule\_aware    & 78/163 (47.9\%) & 75/163 (46.0\%) & 3 \\
\bottomrule
\end{tabular}
\captionof{table}{Bounded-completeness supported rate, raw vs.\ worker-loss-corrected, as a pure rollup
of the \texttt{\_wlc} descriptive CSV.}
\label{tab:p14-comp}
\end{center}

The area raw number, 71/163 (43.6\%), reproduces Chapter~10's headline exactly. The corrected number,
67/163 (41.1\%), lands within one sample of Chapter~10's own estimate of 68/163 (41.7\%) --- and where
Chapter~10 spent 159 fresh GPU reruns to produce that estimate, this reproduces it with \textbf{zero}
reruns, reading columns the descriptive metric now logs by default. The one-sample difference is
explained, not hand-waved: Chapter~10's estimate corrected only top-1 worker loss, whereas the baked-in
verdict discounts worker-loss flips at both top-1 and top-2, so it catches one additional artifact.

\section{A Subtlety Worth Recording}

The correction is deliberately conservative about \emph{which} flips it discounts. A flip is removed
only when the answer changed \emph{and} the worker was lost on that exact rerun --- not whenever a
worker was merely lost. A worker can vanish without the answer changing (the scene-level fallback
happens to return the same answer), and that case is correctly \emph{not} counted as a worker-loss
flip, because nothing about the reported metric was corrupted. Conversely, a genuine flip where the
worker is still detected is never discounted. This keeps ``genuine'' an honest lower bound on true
sensitivity rather than an over-aggressive scrub that would hide real signal.

\section{Standard vs. Non-Standard Choices}

\textbf{Standard.} Reporting a corrected metric alongside a raw one, with the correction's mechanism
stated, is ordinary rigor once a confound is known.

\textbf{Non-standard, and the real contribution.} Two things. First, the correction is promoted from a
one-off, expensive audit (Chapter~10's 159 reruns) into standing metric output computed for free from
data the metric now logs --- the artifact is measured every run, not estimated once. Second, and in
keeping with the whole phase, the raw metric is preserved bit-for-bit behind the default flag, so the
pilot's published raw numbers remain the reproducible baseline against which the genuine numbers are
read; the correction adds a column, it never silently rewrites the number a prior chapter reported.

\section{Implications}

With Phases~1.1 and 1.4 together, the masking and perturbation metrics now both test the right region
\emph{and} report a worker-loss-corrected value, so their aggregates can finally be quoted as measuring
the safety concept rather than the proxy's fallback plumbing. The robustness result sharpens a specific
Phase~2 task: occlusion's inflation is worker loss, which is why targeted and random occlusion must be
separated and reported apart. And completeness is now cheap enough to recompute on every scale run,
because its worker-loss correction no longer carries a 159-rerun tax.
